% SPDX-License-Identifier: Apache-2.0
%
% The cover: what documents TxHDL has, and which to read for what.
% Built by //docs:cover.
\documentclass[journal,onecolumn,9pt]{IEEEtran}

% One column: 80 columns fit at the body size, so nothing is reduced.
\newcommand{\listingsize}{\normalsize}
% SPDX-License-Identifier: Apache-2.0
%
% The house preamble, shared by every document in this directory. Loaded
% after \documentclass, before \title. Keeping it in one file is what
% stops two articles drifting apart in font, listing style or figure
% style.
% Computer Modern. IEEEtran selects Times (ptm) by default, so the face has
% to be chosen explicitly.
%
% Latin Modern is the Computer Modern variety used here, and the choice is
% forced rather than aesthetic. Under T1 encoding, plain `cmr` has no Type 1
% outlines unless cm-super is installed, and it is not in the pinned
% distribution, so pdflatex falls back to EC bitmaps and every glyph in the
% PDF becomes a Type 3 font. That was measured: the first build produced a
% document with 10 Type 3 fonts and no Type 1. Latin Modern is Computer
% Modern redrawn with a full T1 repertoire and real .pfb outlines, so the
% same design comes out scalable.
\usepackage[T1]{fontenc}
\usepackage[utf8]{inputenc}
\usepackage{lmodern}
% microtype is not in the pinned TeX distribution. emergencystretch alone
% keeps the two-column measure from producing overfull lines.
\setlength{\emergencystretch}{3em}

\usepackage{amsmath}
\usepackage{amssymb}
\usepackage{amsthm}
\usepackage{booktabs}
\usepackage{array}
% enumitem is not in the pinned TeX distribution; the list spacing below
% does what \setlist would have done.
\usepackage{float}
\usepackage{xcolor}
\usepackage{listings}
\usepackage{tikz}
\usetikzlibrary{arrows.meta,positioning,fit,backgrounds,calc}
\usepackage[hidelinks,bookmarks=true]{hyperref}

% Numbered, definition-style box for a rule the reader is meant to apply.
\theoremstyle{definition}
\newtheorem{rulebox}{Rule}
\newtheorem{example}{Example}

% Unnumbered asides.
\theoremstyle{remark}
\newtheorem*{tip}{Tip}
\newtheorem*{pitfall}{Pitfall}

% Tighter lists, in place of enumitem's \setlist.
\setlength{\topsep}{2pt}
\setlength{\itemsep}{1pt}
\setlength{\parsep}{0pt}

\newcommand{\code}[1]{\texttt{#1}}
\newcommand{\term}[1]{\emph{#1}}

% Syntax highlighting. Four roles, distinguishable in colour and still
% distinguishable in greyscale, because an IEEE article gets printed: the
% keyword is the only bold role, the comment is the only italic one, and the
% two remaining roles differ in lightness as well as in hue.
\definecolor{kwblue}{RGB}{20,60,150}
\definecolor{cmtgreen}{RGB}{90,120,90}
\definecolor{strred}{RGB}{150,50,40}
\definecolor{numgray}{gray}{0.45}
\definecolor{framegray}{gray}{0.70}
\definecolor{bgshade}{gray}{0.97}
% A darker fill for figure cells. The listing background has to stay very
% light to keep code readable; a figure cell has to be clearly shaded or the
% distinction it draws is invisible in print.
\definecolor{cellshade}{gray}{0.90}

% The listing size is the document's choice, because it depends on the
% column: a two-column article sets `\listingsize` to 6pt and keeps its
% listings within 70 columns, a one-column document keeps the body size
% and includes source files at 80. Lines are never broken; a line that does not fit
% overflows the frame, which is visible, rather than wrapping, which is
% not.
\providecommand{\listingsize}{\scriptsize}

% `'` is deliberately not a string delimiter: the language uses it for loop
% labels, as Rust does, so treating it as a quote would swallow the rest of
% the line.
\lstdefinestyle{txhdl}{
  basicstyle=\ttfamily\listingsize,
  keywordstyle=\bfseries\color{kwblue},
  commentstyle=\itshape\color{cmtgreen},
  stringstyle=\color{strred},
  numberstyle=\tiny\color{numgray},
  morekeywords={mod,use,pub,const,type,struct,enum,trait,impl,fn,async,let,mut,
    match,if,else,loop,while,for,in,break,continue,return,as,await,
    module,bus,channel,transaction,protocol,pipeline,flow,seq,config,
    port,stage,pipe,instance,connect,bind,tag,domain,blackbox,foreign,
    property,role,signal,handler,true,false,
    sequence,interface,modport,implements,package,import,var,wire,func,proc,
    case,default,next,fork,branch,join,state,fsm},
  morecomment=[l]{//},
  morestring=[b]",
  showstringspaces=false,
  columns=fullflexible,
  keepspaces=true,
  breaklines=false,
  % Framed, so a listing is bounded rather than floating in the column.
  frame=single,
  rulecolor=\color{framegray},
  framesep=3pt,
  backgroundcolor=\color{bgshade},
  xleftmargin=3pt,
  xrightmargin=3pt,
  aboveskip=6pt,
  belowskip=6pt,
  % Every listing is numbered and captioned, so it can be referred to by
  % number. `caption` without `float` numbers the listing and leaves it
  % where it was written, which is what a listing in running prose needs.
  captionpos=b,
  abovecaptionskip=2pt,
  belowcaptionskip=6pt,
}

% Figures drawn as text. Framed the same way, and with the highlighting
% switched off, because a box-and-arrow diagram is not code and half its
% words turning blue reads as an accident. Setting a style adds keys rather
% than clearing the ones already set, so the three roles are neutralised by
% hand rather than by leaving them out.
\lstdefinestyle{ascii}{
  basicstyle=\ttfamily\listingsize,
  keywordstyle=\ttfamily,
  commentstyle=\ttfamily,
  stringstyle=\ttfamily,
  columns=fullflexible,
  keepspaces=true,
  frame=single,
  rulecolor=\color{framegray},
  framesep=3pt,
  backgroundcolor=\color{bgshade},
  xleftmargin=3pt,
  xrightmargin=3pt,
  aboveskip=6pt,
  belowskip=6pt,
}
\lstset{style=txhdl}

% House TikZ styles. `shorten` lives on the arrow styles rather than on each
% arrow, so every arrow in the document gets the gap, including ones added
% later. TikZ clips a path at a node's border, so without it an arrowhead
% butts against the box with no gap at all. Keep `node distance` well above
% twice the shorten, or the arrow shrinks to nothing.
\tikzset{
  box/.style={draw,rounded corners=2pt,align=center,inner sep=4pt,
              font=\footnotesize},
  boxf/.style={box,fill=bgshade},
  lbl/.style={font=\scriptsize\itshape,align=center},
  fl/.style={-{Stealth[length=4pt]},shorten >=2.5pt,shorten <=2.5pt},
  fld/.style={fl,dashed},
}



\InputIfFileExists{buildstamp.tex}{}{}
\providecommand{\buildstamp}{Draft build (outside the Bazel flow).}

\title{TxHDL Documents}

\author{Filip Filmar and Dragiša Janković%
\thanks{Both authors are independent. Filip Filmar is the
corresponding author, at f@hdlfactory.com; Dragiša Janković is
reached at linkedin.com/in/dragisa-jankovic.}%
\thanks{This is the cover document of TxHDL. It names every document
the project produces, says what each is for, and says where each is
built. When a document is added, it is added here.}%
\thanks{The authors used a large language model, Claude, as an
assistant in exploring the concepts and in writing and constructing
the documents and the programs; every commit of the repository says
so and carries the prompts that led to it, verbatim.}%
\thanks{\buildstamp}}

% The table of documents fills a page of its own: set it at the top of
% that page rather than at its middle.
\makeatletter
\setlength{\@fptop}{0pt}
\makeatother

\markboth{TxHDL Documents}{Filmar and Janković: TxHDL Documents}

\begin{document}
\maketitle

\section{The Language and Its Documents}

TxHDL is a hardware description language embedded in Rust. It began as
two languages, LHDL and TxHDL, developed independently for the same
project; they were merged into one with a Rust-shaped syntax, and the
merged language was then reduced by deleting every construct Rust
already provides, until what remained was a library. The documents
follow that history, and a reader who wants the current state of the
language reads the last three; a reader who wants to start writing
reads the tutorial.

Every document is built by Bazel from the project repository,
\code{HDL/txhdl}, against a hermetic \LaTeX{} toolchain, so
\code{bazel build //docs/...} on any machine produces all of them.
The release workflow publishes every PDF nightly, and on request.

\section{The Documents}

Table~\ref{tab:documents} names every document, in the order the
combined PDF binds them, with the target that builds each.

\begin{table}[t]
\caption{Every document, what it is for, and where it is built.}
\label{tab:documents}
\centering
\begin{tabular}{@{}p{46mm}p{92mm}l@{}}
\toprule
Document & What it is for & Target \\
\midrule
TxHDL: Unifying a Dataflow and a Transaction HDL
  & The merge. What each source language contributed, the four conflicts
    of meaning and how each was resolved, and the merged language as it
    stood before the embedding. Historical now, and the record of why
    the language has the shape it has.
  & \code{//docs:article} \\
\addlinespace
Deleting a Language: Embedding TxHDL in Rust
  & The language as it is. The rule that every construct Rust provides
    is deleted, what is left, and every claim about Rust compiled as a
    probe. The exposition; read this first.
  & \code{//docs:embedding} \\
\addlinespace
The TxHDL Runtime
  & The library that is the language, included verbatim from the tree,
    with the model of time its executor keeps stated in prose. For
    somebody changing the runtime, or asking what a construct does
    inside.
  & \code{//docs:runtime} \\
\addlinespace
TxHDL by Example
  & One compiled example per concept, from values to an asynchronous
    FIFO between two clocks, each with the output the build produced
    when it ran it. The runtime appears only through its public
    interface. For somebody writing a design.
  & \code{//docs:examples} \\
\addlinespace
Vreteno: An RV32IM Core in TxHDL
  & The first large design: a RISC-V core written in the language,
    checked every cycle against a reference model of the instruction
    set, with its programs and its waveform, and run on an Artix-7
    board, where it said what its simulation said. For somebody
    asking whether the language scales, how a design of this size is
    tested, and whether it has been proven on hardware: it has.
  & \code{//docs:vreteno} \\
\addlinespace
A Tapeout Prep in TxHDL: Vreteno on an Open 45\,nm Node
  & The same core after the FPGA: mapped onto an open 45\,nm standard
    cell library and floorplanned, placed, routed and timed with
    OpenROAD, with the die it came to, the timing it met, and a list of
    everything a real tapeout would still want. For somebody asking
    what happens to a design of this language on the way to silicon.
  & \code{//docs:tapeout} \\
\addlinespace
TxHDL: A Hardware Description Language That Is a Rust Library
  & The paper. The system as it is and what it does, with diagrams:
    the language, the model of time, traces, the lowering and its
    checks, and the Vreteno core through synthesis. No history. For
    somebody meeting TxHDL for the first time.
  & \code{//docs:paper} \\
\addlinespace
TxHDL, Step by Step
  & The tutorial. Six steps on six examples the build runs: a counter,
    the model of time, the three ways of choosing, channels between
    units, the waveform, and the netlist with its checks; then how to
    add an example of your own. For somebody who has written some Rust
    and no hardware.
  & \code{//docs:tutorial} \\
\addlinespace
TxHDL from Zero: A Blinky in Bazel
  & The tutorial for an empty directory: nothing installed, no copy
    of the tree, and four files that build, run and lower a blinky
    with Bazel; the workspace is in the tree and built as a check.
  & \code{//docs:zero} \\
\addlinespace
A Reservation Station in TxHDL
  & The tutorial on a part: what a reservation station is, its rule
    read cycle by cycle on the example's run, how a station of any
    count of inputs is used, with a FIFO behind it, and what its
    netlist is.
  & \code{//docs:station} \\
\addlinespace
AXI in TxHDL
  & The tutorial on a bus: what an AXI link is, the two verbs a host
    client has and the one a peripheral client has, accepting many
    transactions and answering them as the work finishes, the two
    trackers and their netlists, and what routing will attach to.
  & \code{//docs:axi} \\
\addlinespace
A Minimal GPU in TxHDL
  & The first design on the bus: a rasteriser that reads a display
    list and writes pixels into memory over AXI, the framebuffer
    behind the other end of the link, the picture it drew, and the
    netlists of both.
  & \code{//docs:gpu} \\
\addlinespace
A Network on Chip in TxHDL
  & Several hosts and several peripherals, on a lattice: a node with a
    link to each neighbour and an exit, two virtual channels,
    dimension-order routing, and the bridge with AXI on one side and
    packets on the other.
  & \code{//docs:noc} \\
\addlinespace
TxHDL Documents
  & This document.
  & \code{//docs:cover} \\
\addlinespace
TxHDL: What Has Been Built
  & Two pages on the whole of it: the rule the language was built on,
    the parts, the core, the GPU, how all of it is checked, and what
    is not done. For a reader deciding whether to read the rest.
  & \code{//docs:showcase} \\
\addlinespace
TxHDL cheat sheet
  & One page: the rule, the model of time, the vocabulary, a unit with
    its netlist and its waveform, and how to build. Also a PNG.
  & \begin{tabular}[t]{@{}l@{}}\code{//docs:cheatsheet}\\\code{//docs:cheatsheet\_png}\end{tabular} \\
\addlinespace
TxHDL by the Numbers
  & The repository in numbers: lines per topic, per language and per
    file, measured once, on September 13, 2026, and typed in. Not kept
    up to date; the script that measured is in the tree.
  & \code{//docs:stats} \\
\addlinespace
The Dynamics of a Project
  & How the work went rather than what it is: the timeline measured
    from the repository's own history, a Gantt chart of the bands of
    work with what each needed before it could begin, and what a
    repository cannot say.
  & \code{//docs:dynamics} \\
\addlinespace
All of the above, in this order
  & One file, \code{txhdl.pdf}, with a bookmark per document, for
    reading or printing the whole.
  & \code{//docs:all} \\
\bottomrule
\end{tabular}
\end{table}

\section{The Working Notes}

Under \code{docs/} in the repository, in Markdown, are the notes the
articles were written from. They are not typeset, and they are kept
because a decision is worth nothing if the reasons for it are lost.

\begin{itemize}
\item \code{unification-analysis.md}: what each source language is, what
  conflicts, what each contributes, and every defect found in the two
  sources.
\item \code{syntax-decisions.md}: the Rust-shaped surface syntax of the
  merged language, one decision per conflict.
\item \code{rust-embedding.md}: the working notes of the embedding, with
  every probe result in full.
\item \code{document-build-plan.md}: how the documents are built, what
  was tried before, and why the \TeX{} packages are vendored.
\end{itemize}

\section{How the Documents Are Kept True}

Three rules keep the documents and the code from drifting apart, and
they are stated in the repository's \code{AGENTS.md} so that every
change is held to them.

A concept lives in three places or it does not exist: the runtime, a
compiled example, and these documents, in the same change. A listing
of a source file is included from the tree at build time, never typed
into a document, because a typed copy is the one copy nothing checks.
And a claim about what Rust accepts is a probe that was compiled before
it was written down, kept even when the answer was no, because the
error message is the evidence.

\end{document}
